% Table for Section 5.2.2 — one preference pair of each of the two contexts of fig_5_7_context_ranking.
% Generated by result_5_2.ipynb (EXAMPLE_PAIRS gives the pair index per context). Posterior order: Focus, Generic, Probing, Telling.
\begin{table}[!t]
    \centering
    \footnotesize
    \setlength{\tabcolsep}{3pt}
    \renewcommand{\arraystretch}{1.15}
    \captionsetup{skip=4pt}
    \begin{tabularx}{\linewidth}{@{}l>{\raggedright\arraybackslash}Xlr@{}}
        \toprule
        Side & Tutor response (generation model) & \mbox{$p_{\theta}(m \mid r)$} & \mbox{$s(r, e)$} \\
        \midrule
        \multicolumn{4}{@{}l}{\textit{Numerical calculation error} (context 609), $u(m, e)$ = (0.90, 0, 0, 0)} \\
        Chosen & let's look closely at the bridge trolls: the problem says "6 less than four times that number," which means we need to multiply first before subtracting 6. can you try calculating the bridge trolls again using that rule? (gemini-3.5-flash-lite) & \mbox{(0.93, 0.00, 0.05, 0.02)} & 0.83 \\
        Rejected & “6 less than four times 6” means \(4 \times 6 - 6 = 18\) trolls under the bridge, not \(4 \times 6 + 6\). half of 18 is 9, so the total is \(6 + 18 + 9 = 33\) trolls. (gpt-5.6-terra) & \mbox{(0.26, 0.00, 0.07, 0.67)} & 0.24 \\
        \midrule
        \multicolumn{4}{@{}l}{\textit{Missing or incorrect factual knowledge} (context 150), $u(m, e)$ = (0.28, 0, 0.47, 1.67)} \\
        Chosen & the half-off offer must be used on the cheaper original pair, so the \$40 shoes cost \$20 and the \$60 shoes stay \$60, for a subtotal of \$80. a fourth of \$80 is \$20, so james pays \$80 - \$20 = \textbf{\$60}. (gpt-5.6-terra) & \mbox{(0.09, 0.01, 0.03, 0.86)} & 1.48 \\
        Rejected & you correctly found the prices of both pairs of shoes, but let's re-read carefully which pair is supposed to be "half off." which of the two prices, \$40 or \$60, is the cheaper one that gets the discount? (gemini-3.5-flash-lite) & \mbox{(0.88, 0.00, 0.11, 0.01)} & 0.31 \\
        \bottomrule
    \end{tabularx}
    \caption{Second preference pair ($j = 2$) of the two dialogue contexts in Figure~\ref{fig:context_ranking}. The dialogue context (student error type) and the corresponding utility $u(m, e)$ are shown above each pair. Posterior and utility values are listed in the order \textit{Focus}, \textit{Generic}, \textit{Probing}, \textit{Telling}; the score is their inner product.}
    \label{tab:example_pairs}
\end{table}
